\subsection{The one-component cup input from Ohta's isomorphism}
\label{subsec:ohta-component-gate}

We specify the published isomorphism used here, including its recently
completed density input.  Ohta's Part~II, Theorem~(6.2.2), p.~409, in \cite{Ohta2020} proves the isomorphism given
by the upper-right coordinate of the Eisenstein Galois representation.
Its use of CM-point density in Part~I, p.~381, follows from Ohta's \cite[Theorem~(2.3.4), author pp.~15--16]{OhtaDensity}.  This match can be checked directly: the shifts
$\delta_{\mathfrak q}=\{\mathfrak q\}_\Gamma$ chosen on p.~378 of the
2020 paper lie in $\ker(\mathrm{Cl}_\infty\to\mathrm{Cl}_{n_1})$ and
are distinct modulo $\mathrm{Cl}^{\rm alg}$ by its
Proposition~(7.2.4)(iv).  They, together with the unbounded conductor
subsequence $n(\epsilon)$ and the classes in
$\ker(\mathrm{Cl}_{n(\epsilon)}\to\mathrm{Cl}_{n_1})$, are exactly the
data of the 2026 theorem with initial level $n_1$.  Forgetting the
finite level structures maps the former CM tuples onto the latter.
Finiteness and irreducibility of the modular cover then transfer the
density back to the source used in 2020.  Thus the required mod-$p$
nonvanishing input is a theorem.

Here is the action dictionary.  For an even nonexceptional $e$, set
\[
 \chi=\omega^{1-e},\qquad
 \theta_{\mathrm{WWE}}=\omega^{2-e},\qquad
 \theta_{\mathrm O}=\omega^{e-2}.
\]
We define the modular target separately from the arithmetic module
$M=H^{\mathrm{ab}}$.  In the modular-curve and dual-Hecke convention of
\cite[\S3.1]{WWEPseudo}, put $\theta=\theta_{\mathrm{WWE}}$ and
\[
 H'_{\mathrm{mod},\theta}
 =\varprojlim_r H^1_{\mathrm{et}}
   (X_1(p^r)_{\overline{\mathbb Q}},\mathbb Z_p)^{\mathrm{ord}}_\theta,
 \qquad
 H_{\mathrm{mod},\theta}
 =H'_{\mathrm{mod},\theta}\otimes_{\mathfrak h'_\theta}\mathfrak h_\theta.
\]
Here $X_1(p^r)$ is the compactified modular curve classifying
elliptic curves with a point of exact order $p^r$, and the ordinary
summand is the summand on which the dual $U_p$ operator is invertible.
The transitions are corestrictions, the subscript is the diamond
eigenspace, and $\mathfrak h_\theta$ is the localization of the algebra
$\mathfrak h'_\theta$ of dual Hecke operators at its Eisenstein maximal
ideal.  Let $I_\theta\subset\mathfrak h_\theta$ be the cuspidal Eisenstein
ideal, the image of the annihilator of the ordinary Eisenstein series
in the full Hecke algebra.  Set
\[
 M_{\mathrm{mod}}
 =H^-_{\mathrm{mod},\theta}/I_\theta H^-_{\mathrm{mod},\theta},
\]
where the minus sign denotes complex conjugation.  The coefficient
$\Lambda$ acts by $\gamma\mapsto\langle\kappa(\gamma)\rangle$.

Ohta's class character is
$(\theta_{\mathrm O}\omega)^{-1}=\chi$.  His tame level is $1$,
and the exponent $e-2$ is even, so that $-1$ belongs to its kernel,
as required in Theorem~(6.2.2).  The exceptional indices $0$ and $2$
have $T_0=T_2=0$ and need no cup-symbol argument.  Ohta's $X^-$ is a
splitting of the ordinary quotient $X/X^+$, with $X^+=X^{I_p}$
in his (6.1.14).  The comparison with the above minus space uses the
natural isomorphism $H^-_{\mathrm{mod},\theta}\to H_{\mathrm{quo},\theta}$
of \cite[Theorem~3.4.1 and Corollary~3.4.2]{WWEPseudo}, followed by the
modular-curve and diamond-action transport to that ordinary quotient.
In Ohta's (6.1.26), an element $\gamma$ with $\kappa(\gamma)=u$ acts on
$(X^-/I^*X^-)^\dagger$ as $u^{-1}\iota(u)^{-1}$.  After the Tate twist $(1)$,
this becomes the inverse diamond action $\iota(u)^{-1}$; its character
is $\theta_{\mathrm{WWE}}$.  The extension-coordinate map is
\[
 \sigma\longmapsto \sigma e^+-e^+\pmod{I^*X}.
\]
Sharifi's \cite[pp.~3 and 13--14]{SharifiProgram} explicitly uses the inverse-Galois convention and records the
Ohta isomorphism for this $\Upsilon$.  Equivalently, this is the
upper-right-coordinate definition in \cite[\S8.1]{WWEPseudo}:
evaluate the extension coordinate in
$\operatorname{Hom}_{\mathfrak h_\theta/I_\theta}
 (H^+_{\mathrm{mod},\theta}/I_\theta H^+_{\mathrm{mod},\theta},
 M_{\mathrm{mod}})$ on the distinguished element corresponding to
$1$ under $\mathfrak h_\theta/I_\theta\simeq
H^+_{\mathrm{mod},\theta}/I_\theta H^+_{\mathrm{mod},\theta}$.
Its restriction to the cyclotomic kernel is a homomorphism, and
ordinarity makes it factor through $X_\chi(1)$.
We therefore have the isomorphism
\[
 \Upsilon: V=X^{[1-e]}(1)\xrightarrow{\ \sim\ }M_{\mathrm{mod}}.
\]

Suppose now that $X^{[e]}$ is finite, so that the previously established
adjoint calculation makes $V[1/p]$ cyclic.  At a characteristic prime
$f$, let $R$ be the corresponding discrete valuation ring and write
$V_R\simeq R/(f^r)$.  Transport this generator to
$(M_{\mathrm{mod}})_R$ by $\Upsilon$.
The endomorphism $\Upsilon\varpi$ is multiplication by an element $a$,
and the Fukaya--Kato identity in
\cite[Theorem~5.21]{SharifiFK} gives
\[
 (a-1)\xi'=0\quad\hbox{in }R/(f^r).
\]
Writing $\xi=f^r g$ with $g\in R^\times$, separability in
characteristic zero gives $v_f(\xi')=r-1$.  Hence $a-1\in(f)$,
so $a$ is a unit.  Thus $\varpi_R$ is an isomorphism at every
characteristic prime, with no restriction on the multiplicity $r$.
It follows that the rational cup-symbol map is surjective on all of
$V[1/p]$. Both integral modules are finitely generated over $\Lambda$
with zero $\mu$-invariant, so the cokernel of the integral map
is a finite $\mathbb Z_p$-module. One power of $p$ kills it.
Reduction along the norm-compatible finite-level symbol maps
therefore gives the uniform integral bound used below.

Thus $X^{[e]}$ finite supplies cyclicity and the full rational
cup image in the twice-twisted companion.  The arithmetic character
transport is the displayed map $\theta_{\mathrm O}\mapsto
\theta_{\mathrm{WWE}}$ together with the Tate twist.

The cyclicity step is used in the direction displayed in the
Fukaya--Kato identity. For example, in a discrete valuation ring
with uniformizer $t$ and residue characteristic zero, take
$V=R/(t^2)$, $M_{\mathrm{mod}}=(R/t)^2$ and both maps zero. Then
$t(\Upsilon\varpi-1_{M_{\mathrm{mod}}})=0$, while $tV\ne0$.
In the argument above, Ohta's isomorphism identifies $M_{\mathrm{mod}}$ with the
cyclic module $V$ and turns the identity into multiplication by
$a$; that is the step which excludes this example.
